\originalsection{微观漂移与电阻率的条件}\label{tb:drude}
在简单碰撞模型中, 电子每次碰撞后失去先前定向运动记忆, 平均自由时间为$\tau_c$, 电场提供加速度$-e\ee/m_e$. 平均漂移$\vct v_d=-e\tau_c\ee/m_e$, 因而
\[\jj=(-e)n\vct v_d=\frac{ne^2\tau_c}{m_e}\ee,\quad
\sigma_c=\frac{ne^2\tau_c}{m_e}.\]
随机热运动速度通常远大于漂移速度, 但随机方向平均相消. 这是Drude式的现象模型, 不意味着全部金属、半导体、超导体的电阻可由一条经典碰撞公式精确解释. 材料、温度和频率会改变响应.

在有限温区常以$\rho_r(T)\simeq\rho_r(T_0)[1+\alpha(T-T_0)]$作线性近似. $\alpha$单位为$\mathrm{K^{-1}}$, 金属通常正值, 载流子浓度变化显著的半导体可能负值. 这不是任意温度外推的定律. 电阻$R$还依赖几何, 不应把温度依赖的电阻率和一段导线的电阻互换.

均匀圆线内每单位体积吸收功率为$\jj\cdot\ee=\sigma_cE^2$. 体积积分给$I^2R$. 沿导线需要电场来维持碰撞中的漂移, 所以有电流的金属材料内部一般不是静电平衡的零场导体. 表面微小电荷分布建立导线内的场, 这与前章静止等势导体的条件不同.

\begin{example}
\textbf{原6.5.1: 早期跨海电缆.} 
长$3000\,\mathrm{km}$的缆芯由七根直径$0.73\,\mathrm{mm}$的铜线并联组成, 电阻率给为$3.0\times10^{-6}\,\Omega\mathrm{cm}$. 忽略接头和交流效应, 求电阻.
\end{example}
\begin{solution}
先换单位$\rho_r=3.0\times10^{-8}\,\Omega\mathrm m$. 总截面$A=7\pi(0.73\times10^{-3}/2)^2\simeq2.93\times10^{-6}\,\mathrm{m^2}$, 故$R=\rho_rL/A\simeq3.07\times10^4\,\Omega$. 七根并联改变面积, 长度仍是整段电缆长度. 值是题给电阻率模型的结果, 不能用现代室温铜常数替换后仍称复现原题.
\end{solution}

\begin{example}
\textbf{原6.5.2: 两导电材料界面的电荷.} 
两材料电导率为$\sigma_1,\sigma_2$, 截面积同为$A$, 稳态电流$I$垂直界面从1流向2. 原题采用$\eps$的法向场跃变关系. 求界面自由电荷.
\end{example}
\begin{solution}
无持续界面积累要求$J_1=J_2=I/A$. 欧姆关系给$E_i=I/(A\sigma_i)$, 所以
\[Q_s=A\eps(E_2-E_1)=\eps I\left(\frac1{\sigma_2}-\frac1{\sigma_1}\right).\]
较差导体需要较强电场, 界面电荷提供这个跃变. 若两材料还有不可忽略的不同介电率, 自由电荷应从$D_{n,2}-D_{n,1}$算, 原来的$\eps$式不应无条件照搬.
\end{solution}

\begin{example}
\textbf{原6.5.3: 海水中的双向离子电流.} 
海水电阻率$25\,\Omega\mathrm{cm}$, 钠离子、氯离子各$3.0\times10^{20}\,\mathrm{cm^{-3}}$. 长$2\,\mathrm m$水管两端电压$12\,\mathrm V$, 假定两类一价离子漂移速率相同, 求平均漂移速率.
\end{example}
\begin{solution}
$E=6\,\mathrm{V/m}$, $\rho_r=0.25\,\Omega\mathrm m$, $J=E/\rho_r=24\,\mathrm{A/m^2}$. 正离子顺场、负离子逆场, 但约定电流都顺场, 所以$J=e(n_++n_-)v_d$而非两者相减. 代$n_i=3.0\times10^{26}\,\mathrm{m^{-3}}$得$v_d\simeq2.50\times10^{-7}\,\mathrm{m/s}=2.50\times10^{-5}\,\mathrm{cm/s}$. 若迁移率不同, 题给数据只确定$n_+v_++n_-v_-$, 不能各自唯一求出两个速度.
\end{solution}

\begin{example}
\textbf{原6.5.4、6.5.5: 电阻几何.} 
(a)高$h$的截锥电阻, 两端半径$b,a$, 电流沿轴. (b)长$L$、内外半径$a,b$的中空圆柱, 分别沿轴和沿径向导电. 求电阻.
\end{example}
\begin{solution}
(a)采用每一截面近似等势的一维缓变模型, $r(x)=b+(a-b)x/h$, 小片串联:
\[R=\int_0^h\frac{\rho_r\dd x}{\pi r(x)^2}=\frac{\rho_r h}{\pi ab}.\]
这近似忽略横向场, 剧烈变化的实际锥体需解三维边界问题. $a=b$还原圆柱.

(b)轴向截面为$\pi(b^2-a^2)$, 所以$R_z=\rho_rL/[\pi(b^2-a^2)]$. 径向两电极覆盖内外圆柱面、忽略端部时, 小壳导电面积为$2\pi rL$, 得
\[R_r=\int_a^b\frac{\rho_r\dd r}{2\pi rL}=\frac{\rho_r}{2\pi L}\ln\frac ba.\]
径向电流不是穿过环形截面, 用轴向面积代替会得错误尺度.
\end{solution}

\originalsection{回路独立性、电源功率与立方体电阻}\label{tb:kirchhoff-examples}
回路方程本质是电压降与电源升势的能量账. 每个含电阻支路的电流方向可以先任取, 相应电压降$IR$也保留代数符号. 遍历方向相反时整项反号. 网络有$N$个节点、$B$条支路且连通时, 独立节点方程最多$N-1$个, 独立回路最多$B-N+1$个. 这是避免重复方程的实用数数法.

\begin{example}
\textbf{原7.9.1: 求给定等效电阻的元件值.} 
输入电阻为$R_1$串联$R_1$与$(R_0+R_1)$的并联, 要求等效值恰为$R_0>0$. 求正的$R_1$.
\end{example}
\begin{solution}
$R_0=R_1+R_1(R_0+R_1)/(R_0+2R_1)$. 乘分母整理为$3R_1^2=R_0^2$, 因而$R_1=R_0/\sqrt3$. 负根不是被动电阻的可接受值.
\end{solution}

\begin{example}
\textbf{原7.9.2、7.9.4: 最大功率与串并联功率.} 
电动势$\mathcal E$、内阻$r>0$的电池接负载$R$. 求最大负载功率时的$R$和效率. 另有两个电阻接固定理想电压$V$, 比较串、并联时各自及总功率.
\end{example}
\begin{solution}
$I=\mathcal E/(R+r)$, $P_R=\mathcal E^2R/(R+r)^2$, 导数为$\mathcal E^2(r-R)/(R+r)^3$, 所以$R=r$为最大值$\mathcal E^2/(4r)$. 电源总化学功率$\mathcal EI$, 效率$P_R/(\mathcal EI)=R/(R+r)$, 匹配时仅一半. 高效率与最大输出功率是不同目标.

并联各自$P_i=V^2/R_i$, 总功率$V^2(1/R_1+1/R_2)$. 串联电流$V/(R_1+R_2)$, 各自$P_i=V^2R_i/(R_1+R_2)^2$, 总量$V^2/(R_1+R_2)$. 大电阻在串联中发热更多, 在并联中更少. 两种配置中各电阻功率之和都等于电源的$VI$. 固定同一理想端电压下,$1/R_1+1/R_2>1/(R_1+R_2)$, 所以并联总功率严格大于串联, 完成原7.9.4(e)比较. 真实灯泡的电阻随温度改变, 亮度也非简单恒定电阻定律的直接定义.
\end{solution}

\begin{example}
\textbf{原7.9.5: 立方体十二条边上的电阻.} \label{ex:resistorcube}
立方体每条棱均为电阻$R$, 在一对相对顶点间接电源. 求等效电阻.
\end{example}
\begin{solution}
从起点到终点的体对角线具有三重旋转对称. 起点相邻的三个顶点等势, 终点相邻的另三个顶点也等势. 于是第一层三条边并联为$R/3$, 中间两等势组由六条边连接为$R/6$, 最后一层又是$R/3$. 三层串联, 总量$5R/6$. 等势是由完整网络与端点的对称性推出, 不能在任意两个空间邻近点之间擅自短接.
\end{solution}

\originalsection{RC开关中的连续量与能量}\label{tb:rc-switch}
\begin{example}
\textbf{原7.9.3: 充电支路被短接后的电流.} 
电池$\mathcal E$经过$R_1$接节点A. A到电池负端的支路依次为电容$C$及$R_2$, 原电容正板直接接A, 还有起初断开的理想开关. 断开很久后闭合开关, 求电容电荷、$R_2$电流及开关电流.
\end{example}
\begin{solution}
闭合前电流为零, 电容电压为$\mathcal E$, $q_0=C\mathcal E$. 初始充电时间常数是$(R_1+R_2)C$. 闭合后A被接到负端, 电容只经过$R_2$放电, 故
\[q=C\mathcal E e^{-t/(R_2C)},\qquad |I_2|=\frac{\mathcal E}{R_2}e^{-t/(R_2C)}.\]
放电电流从电容正板直接进入A, 经闭合开关流向负端, 再经$R_2$回到电容负板, 与原充电支路的方向相反. 电池持续经$R_1$提供$\mathcal E/R_1$, 两者在A汇合, 开关流向负端的电流是
\[I_s=\frac{\mathcal E}{R_1}+\frac{\mathcal E}{R_2}e^{-t/(R_2C)}.\]
原解若用负的支路电流变量, 节点式必须同时保留方向; 不能把负幅值与“流入量相加”混用. 电容电压在$t=0$连续, 所以初始并非无电压的短路.
\end{solution}

初始无电荷的RC充电, 电源供能$\int_0^\infty\mathcal EI\dd t=C\mathcal E^2$. 最终储能$C\mathcal E^2/2$, 电阻热$\int_0^\infty I^2R\dd t=C\mathcal E^2/2$. 这后一半与$R$数值无关, 但$R$决定释放速率. 把$R$直接设为零会使理想充电电流成为冲量, 物理电路由电感、辐射和电源内阻等响应, 不能据理想极限否认能量守恒.

汽车启动时大电流使电池内阻上的压降增大, 端电压下降而灯变暗. 短接具有内阻的电池时端电压可近零, 化学能仍在内阻中耗散. RC电阻改变充电时间却不改变恒压下最终电荷. 这些原7.10概念题都依赖先区分电源内部和外部元件.

\originalsection{分压的负载条件与不能串并联化简的网络}\label{tb:two-source}
两个串联电阻$R_1,R_2$接理想源$V_s$, 空载输出跨$R_2$时$V_o=V_sR_2/(R_1+R_2)$. 接负载$R_L$以后, 必须以$R_2\parallel R_L$替代$R_2$. 只有$R_L\gg R_2$且其他源内阻可忽略, 原来的空载分压才近似不变. 这是原7.4电阻分压器的实际使用条件.

\begin{example}
\textbf{原7.8: 两个电源的三支路.} 同一左、右节点间, 中支路含$\mathcal E_2,R_1$, 上支路含$R_2,\mathcal E_1$, 下支路只有$R_3$. $\mathcal E_2$正端朝右, $\mathcal E_1$正端朝左. 原图定义$I_1,I_2$从左流向右, $I_3$从右流向左. 求三支路电流.
\end{example}
\begin{center}
\begin{tikzpicture}[>=Stealth]
\draw (0,-1.5)--(0,1.5)--(.8,1.5);\draw (1.8,1.5)--(2.7,1.5);\draw (3,1.5)--(5,1.5)--(5,-1.5);
\draw (.8,1.25) rectangle (1.8,1.75);\node at (1.3,1.5) {$R_2$};
\draw[thick] (2.7,1.12)--(2.7,1.88);\draw (3,1.3)--(3,1.7);\node[above] at (2.85,1.9) {$\mathcal E_1$};\node[below] at (2.6,1.1) {$+$};
\draw (0,0)--(1.1,0);\draw (1.4,0)--(2.7,0);\draw (3.7,0)--(5,0);
\draw (1.1,-.2)--(1.1,.2);\draw[thick] (1.4,-.38)--(1.4,.38);\node[above] at (1.25,.4) {$\mathcal E_2$};\node[below] at (1.5,-.4) {$+$};
\draw (2.7,-.25) rectangle (3.7,.25);\node at (3.2,0) {$R_1$};
\draw (0,-1.5)--(2,-1.5);\draw (3,-1.5)--(5,-1.5);\draw (2,-1.75) rectangle (3,-1.25);\node at (2.5,-1.5) {$R_3$};
\draw[->,structurecolor] (3.6,1.8)--(4.6,1.8) node[midway,above] {$I_2$};
\draw[->,structurecolor] (3.8,.3)--(4.8,.3) node[midway,above] {$I_1$};
\draw[->,structurecolor] (1.7,-1.9)--(.7,-1.9) node[midway,below] {$I_3$};
\end{tikzpicture}
\end{center}
\begin{solution}
记$V=V_{\rm right}-V_{\rm left}$, 三支路关系为
\[I_1=\frac{\mathcal E_2-V}{R_1},\quad I_2=\frac{-\mathcal E_1-V}{R_2},\quad I_3=\frac V{R_3},\qquad I_1+I_2=I_3.\]
代入节点式, 分母$D=R_1R_2+R_1R_3+R_2R_3$, 得
\[V=\frac{R_3(\mathcal E_2R_2-\mathcal E_1R_1)}D,\quad
I_1=\frac{\mathcal E_1R_3+\mathcal E_2R_3+\mathcal E_2R_2}D,\]
\[I_2=-\frac{\mathcal E_1R_1+\mathcal E_1R_3+\mathcal E_2R_3}D,\quad
I_3=\frac{\mathcal E_2R_2-\mathcal E_1R_1}D.\]
对正的两电动势,$I_2<0$表示真实上支路电流反向, 并非电流大小负. 等价的三独立方程是节点守恒和两个回路式$\mathcal E_2-I_1R_1-I_3R_3=0$, $-I_2R_2-\mathcal E_1+I_1R_1-\mathcal E_2=0$. 绕外圈的新式只是这两式相加, 不增加独立信息. 电源总提供功率$\mathcal E_2I_1-\mathcal E_1I_2$等于$\sum R_iI_i^2$, 由支路式乘电流相加即得. 图在本稿横向重排, 实际三个节点支路、电源极性和电流定义保持原件.
\end{solution}

\originalsection{实际仪表的扰动与暂态读图}\label{tb:meters}
测量电阻$R$的电压时, 有限内阻$R_V$电压表与它并联, 网络实际负载变为$R\parallel R_V$. 若源有内阻$r_s$, 读数为$\mathcal E(R\parallel R_V)/[r_s+(R\parallel R_V)]$, 与不接表的$\mathcal ER/(r_s+R)$可不同. 电流表内阻$R_A$串入后, 单回路读数变成$\mathcal E/(r_s+R+R_A)$, 极小$R_A$才近无扰动. 理想源直接固定两端电压的特例可掩盖电压表加载, 因而要明确源条件. 电路图的ground通常是任选零势节点, 物理接地则实际连接大地, 两者不必同义.

原7.5还示范电阻色环. 前两环给有效数字、第三环给十的幂次倍率、第四色环给容差; 黄、橙、黑、金表示$43\times10^0\,\Omega$、$\pm5\%$. 常用数字色码为黑0、棕1、红2、橙3、黄4、绿5、蓝6、紫7、灰8、白9. 色环是制造标记约定, 不决定电阻的实际温度和功率响应; 测量值应连同容差理解.

未充电RC合闸后,$q/(C\mathcal E)=V_C/\mathcal E=1-e^{-t/\tau}$, $IR/\mathcal E=e^{-t/\tau}$. 在$t=0,\tau,\infty$时, 电容比例依次为$0,0.632,1$, 电流比例为$1,0.368,0$. 这使“时间常数”成为可从图或测量辨认的尺度. 放电则电压从初值按同一衰减曲线下降.
\begin{center}
\begin{tikzpicture}
\begin{axis}[width=.8\linewidth,height=4.7cm,xlabel={$t/\tau$},ylabel={归一化值},xmin=0,xmax=5,ymin=0,ymax=1.1,samples=101,legend pos=outer north east]
\addplot[thick,structurecolor,domain=0:5] {1-exp(-x)};\addlegendentry{$q/(C\mathcal E)$}
\addplot[thick,dashed,domain=0:5] {exp(-x)};\addlegendentry{$IR/\mathcal E$}
\end{axis}
\end{tikzpicture}
\end{center}

原7.11.4无限电阻梯引入自相似方法. 以一节串联$R_s$后分出并联$R_p$、其余部分又与全梯相同的模型说明: 全梯输入电阻$X$满足$X=R_s+(R_p\parallel X)$. 整理
\[X^2-R_sX-R_sR_p=0,\quad X=\frac{R_s+\sqrt{R_s^2+4R_sR_p}}2.\]
负根不是被动输入电阻. 有限$n$节截断网络满足同一迭代$X_{n+1}=R_s+R_pX_n/(R_p+X_n)$, 当$X_n\ge R_s>0$时映射导数不超过$[R_p/(R_p+R_s)]^2<1$, 所以终端扰动逐节缩小、极限由正根确定. 自相似方程必须由实际重复单元的连接推出, 不能对任意画得很长的网络猜同一二次式. 此处是该方法的明确迁移模型, 不是未经核对改写原自测题的数值.
